\documentclass[
 aps,
 prd,
 reprint,
 amsmath,
 amssymb,
 nofootinbib,
 superscriptaddress
]{revtex4-2}

\usepackage[T1]{fontenc}
\usepackage{lmodern}
\usepackage{bm}
\usepackage{mathrsfs}
\usepackage{amsthm}
\usepackage{booktabs}
\usepackage{array}
\usepackage{enumitem}
\usepackage{microtype}
\usepackage{hyperref}
\usepackage{orcidlink}
\usepackage[nameinlink,capitalize]{cleveref}

\hypersetup{
  colorlinks=true,
  linkcolor=blue,
  citecolor=blue,
  urlcolor=blue
}

\newtheorem{theorem}{Theorem}[section]
\newtheorem{proposition}[theorem]{Proposition}
\newtheorem{lemma}[theorem]{Lemma}
\newtheorem{corollary}[theorem]{Corollary}
\theoremstyle{definition}
\newtheorem{definition}[theorem]{Definition}
\theoremstyle{remark}
\newtheorem{remark}[theorem]{Remark}

\newcommand{\Sthree}{S^3_R}
\newcommand{\Vol}{\operatorname{Vol}}
\newcommand{\PP}{\mathbb P}
\newcommand{\RR}{\mathbb R}
\newcommand{\cK}{\mathfrak K}
\newcommand{\cKp}{\mathfrak K_+}
\newcommand{\cKm}{\mathfrak K_-}
\newcommand{\cV}{\mathcal V}
\newcommand{\cH}{\mathcal H}
\newcommand{\cL}{\mathcal L}
\newcommand{\Eul}{\mathsf E}
\newcommand{\ad}{\operatorname{ad}}
\newcommand{\curl}{\operatorname{curl}}
\newcommand{\divg}{\operatorname{div}_g}
\newcommand{\gradg}{\operatorname{grad}_g}
\newcommand{\Def}{\operatorname{Def}}
\newcommand{\dd}{\,\mathrm d}
\newcommand{\Id}{\mathrm{Id}}
\newcommand{\rank}{\operatorname{rank}}

\begin{document}

\title{Killing--Beltrami algebra and exact reductions of Navier--Stokes flow on the three-sphere}

\author{Boris~Batenin\,\orcidlink{0009-0000-7027-8583}}
\email{batenin@web.de}
\affiliation{Independent Researcher, Berlin, Germany}

\author{Andrei~Preece\,\orcidlink{0009-0004-5978-795X}}
\email{andre.preece@web.de}
\affiliation{Independent Researcher, Berlin, Germany}

% APS submission-form action still required:
% designate exactly one corresponding author in the APS submission server.

\begin{abstract}
We identify an exact $SO(4)$-equivariant interaction algebra for incompressible
Navier--Stokes flow with Ebin--Marsden viscosity on the round three-sphere.
The central identity determines the polarized interaction of every Killing
field with every signed curl eigenspace: on $E_{k,\sigma}$, the two simple
factors of the Killing algebra act by Lie derivatives with distinct
chirality-dependent coefficients. This intertwining law organizes a
shell-diagonal skew dynamics, a canonical generally nonflat metric connection
over $SO(4)$, exact viscous finite-dimensional reductions, and a rigidity
theorem stating that every finite exact reduction containing the full Killing
algebra is a union of complete signed curl shells. We also obtain a
quadratic-Casimir representation of the Stokes operator and viscous
dissipation in terms of six Killing convective responses. Compact homogeneity
then yields sharp projector norms for finite $SO(4)$-invariant vector sectors,
with sharp spectral concentration consequences. The maximal
torus gives a separate infinite-dimensional pressure-linear sector. These
geometric and spectral statements are intrinsic to the round three-sphere and
are logically independent of the conditional singularity-transfer
construction that motivated the analysis.
\end{abstract}

\maketitle




\section{Introduction and main results}
\label{sec:intro}

Geometric formulations of incompressible fluid dynamics on Riemannian
manifolds make the interaction between viscosity, isometries, and curvature
structural rather than perturbative. We use the deformation-tensor
(Ebin--Marsden) viscosity on the round three-sphere $S^3_R$, following the
geometric framework of Ebin and Marsden and subsequent formulations of
Navier--Stokes flow on Riemannian manifolds
\cite{EbinMarsden1970,ChanCzubakDisconzi2017,SamavakiTuomela2020}.
The signed curl spectrum on the round sphere is classical
\cite{IkedaTaniguchi1978}. The geometric and spectral program from which the
present work grew was developed by Batenin and Preece in a Zenodo preprint
\cite{PreeceBatenin2026}; that work combines an intrinsic theory of exact
spectral closure on $S^3_R$ with a separate conditional transfer of a forced
singular branch.

Two pieces of prior literature are important for positioning the results
below. Samavaki and Tuomela proved that, for the same deformation-tensor
viscosity on a compact Riemannian manifold, the component of an unforced
solution along any Killing field is conserved
\cite{SamavakiTuomela2020}. The bilinear Killing cancellation used below is
the polarized form of that geometric cancellation and is therefore not our
principal novelty. More recently, Heslin and Preston constructed explicit
nonstationary Euler solutions on curved manifolds, including the round
three-sphere \cite{HeslinPreston2026}. Thus neither conservation of Killing
momentum nor the mere existence of exact Killing/eigenfield waves is the
distinguishing point of the present article.

The new organizing mechanism is the exact polarized Killing--Beltrami
interaction. Let
\[
 \mathfrak K=\mathfrak K_+\oplus\mathfrak K_-
 \simeq\mathfrak{su}(2)_+\oplus\mathfrak{su}(2)_-
\]
be the six-dimensional Killing algebra, with
$\mathfrak K_\tau=E_{0,\tau}$, and let
$v\in E_{k,\sigma}$ be a signed curl eigenfield. We prove
\begin{equation}
 \boxed{
 2Q(K_\tau,v)
 =
 \frac{k+2-2\tau\sigma}{k+2}[K_\tau,v].
 }
 \label{eq:intro-intertwining}
\end{equation}
The two simple factors therefore act on the same signed shell with the
different coefficients
\begin{equation}
 \frac{k}{k+2},
 \qquad
 \frac{k+4}{k+2}.
 \label{eq:intro-chiral-coeffs}
\end{equation}
This is not merely an exact-solution ansatz: it is a complete intertwining law
for the polarized convective action of the Killing algebra on every signed
curl shell.

Several results follow from the same identity. First, the Killing action on
the orthogonal complement is shell diagonal and skew-adjoint. It transfers
neither $L^2$ energy nor helicity between signed shells; all inter-shell
nonlinear transfer is carried by the high--high term $Q(v,v)$. Second, on
each $E_{k,\sigma}$ the map $m\mapsto D_{k,\sigma}(m)=2Q(m,\cdot)$ fails in a
controlled way to be a Lie-algebra representation. Its defect is exactly the
curvature of a canonical metric connection on
\[
 SO(4)\times E_{k,\sigma}\longrightarrow SO(4),
\]
and the physical Killing precession is the parallel transport of this
connection. The curvature coefficients decay as $k^{-1}$, but the
representation generators grow like $k$; consequently the natural operator
norm of the curvature does not flatten in the ultraviolet.

Third, the intertwining formula produces finite-dimensional exact
Navier--Stokes reductions
\[
 \mathfrak K_\tau\oplus E_{k,\sigma},
 \qquad
 \mathfrak K\oplus E_{k,\sigma},
\]
with explicit triangular viscous evolution. More importantly, we prove a
converse rigidity theorem: if a finite-dimensional exact reducing space
contains the full Killing algebra, then every nonzero higher spectral
component is a complete signed curl shell.

Fourth, if $T_a^\tau=Q(K_a^\tau,\cdot)$ for an $L^2$-orthonormal Killing
basis, then
\begin{equation}
 \boxed{
 A=-2V_R\sum_{\tau=\pm}\sum_{a=1}^3(T_a^\tau)^2,
 }
 \label{eq:intro-casimir}
\end{equation}
so the Stokes operator and viscous dissipation are represented by the
quadratic Casimir of the six convective Killing responses.

Compact homogeneity supplies the complementary spectral statement
\begin{equation}
 \|\Pi_W\|_{L^2\to L^\infty}
 =
 \sqrt{\frac{N}{3V_R}},
 \qquad
 \|\Pi_W\|_{L^1\to L^\infty}
 =
 \frac{N}{3V_R},
 \label{eq:intro-projectors}
\end{equation}
for every finite-dimensional $SO(4)$-invariant vector sector $W$ of dimension
$N$, with sharp constants.

The exterior-parametrix arithmetic, Mellin shift hierarchy, interaction
picture, and polynomial recursion are not used in the proofs below and are
outside the scope of the present article. Likewise, the conditional forced
singularity transfer is not used to prove any intrinsic result here.


\section{Geometric and spectral setting}
\label{sec:setting}

Let $(\Sthree,g)$ be the round three-sphere of radius $R$ and let
\begin{align*}
 \partial_t u+\nabla_u u-\nu\cL_g u+\gradg p&=f,\\
 \divg u&=0,\\
 \cL_g u&=2\,\operatorname{div}_g\Def_g u .
\end{align*}
Write
\begin{align*}
 A&=-\PP\cL_g,\\
 B(u,v)&=\PP\nabla_u v,\\
 Q(u,v)&=\frac12\{B(u,v)+B(v,u)\}.
\end{align*}
Thus the projected equation is
\begin{equation}
 \partial_tu+\nu Au+Q(u,u)=\PP f.
 \label{eq:NSproj}
\end{equation}

The spectral conventions used here agree with Ref.~\cite{PreeceBatenin2026}; in particular, the signed curl decomposition
\[
 L^2_\sigma(T\Sthree)
 =
 \widehat{\bigoplus}_{k\ge0,\ \sigma=\pm}E_{k,\sigma},
\]
with
\begin{align}
 \curl u&=\mu_{k,\sigma}u,
 & \mu_{k,\sigma}&=\sigma\frac{k+2}{R},
 \label{eq:curlspec}\\
 Au&=\kappa_k u,
 & \kappa_k&=\frac{k(k+4)}{R^2},
 \label{eq:stokespec}\\
 \dim E_{k,\sigma}&=d_k=(k+1)(k+3).
 \label{eq:dim}
\end{align}
Moreover,
\begin{equation}
 E_{k,+}\simeq V_{(k+2)/2}\boxtimes V_{k/2},
 \qquad
 E_{k,-}\simeq V_{k/2}\boxtimes V_{(k+2)/2}
 \label{eq:irreps}
\end{equation}
under
$SO(4)\simeq(SU(2)_+\times SU(2)_-)/\{\pm1\}$.
The $k=0$ shells are the two simple Killing factors:
\[
 \cK=\cKp\oplus\cKm,
 \qquad
 \cK_\tau:=E_{0,\tau},
 \qquad \dim\cK_\tau=3.
\]

For later comparison, Ref.~\cite{PreeceBatenin2026} also records
\begin{equation}
 Q(\cK,\cK)=0,
 \qquad
 Q(E_{k,\sigma},E_{k,\sigma})=0,
 \label{eq:diagonalzero}
\end{equation}
and introduces the complete-shell leakage scalar
\[
 \mathfrak R(S)=
 \sum_{\alpha,\beta\in S}\sum_{\gamma\notin S}
 \left\|\Pi_\gamma
 Q\big|_{E_\alpha\times E_\beta}\right\|_{\mathrm{HS}}^2.
\]
For a supplied finite union
$W_S=\bigoplus_{\alpha\in S}E_\alpha$,
Ref.~\cite{PreeceBatenin2026} proves
\begin{equation}
 W_S\text{ is exact reducing}
 \quad\Longleftrightarrow\quad
 \mathfrak R(S)=0.
 \label{eq:Rcriterion}
\end{equation}

We derive consequences of these intrinsic facts.  No result
below uses the generic-order inner solver, renormalized matching hypothesis,
or annular-cycle stability hypothesis required by the relative singular
transfer in Ref.~\cite{PreeceBatenin2026}.


The new statements use only intrinsic inputs on the round sphere:
the metric, curl and Leray projection for the Killing--Beltrami interaction;
the Stokes spectrum and irreducibility for Killing-saturated classification;
the torus reduction and Jacobi spectrum for the maximal-torus sector; and
compact homogeneity with irreducible isotropy for the projector identities.





\section{The Killing--Beltrami interaction algebra}
\label{sec:KB}

We first isolate the identity that drives the new finite-dimensional
closures.

\begin{lemma}[Polarized vector identity]
For smooth vector fields $u,v$ on an oriented Riemannian three-manifold,
\begin{equation}
 \nabla_u v+\nabla_v u
 =
 \gradg\langle u,v\rangle
 -u\times\curl v-v\times\curl u.
 \label{eq:polarvec}
\end{equation}
If additionally $\divg u=\divg v=0$, then
\begin{equation}
 \curl(u\times v)=-[u,v].
 \label{eq:curlcross}
\end{equation}
\end{lemma}

\begin{proof}
Both identities are tensorial versions of the standard three-dimensional
vector-calculus formulas.  For \eqref{eq:polarvec}, pair both sides with an
arbitrary vector field and use
$\dd(v^\flat)(X,Y)=\langle\nabla_Xv,Y\rangle-
\langle\nabla_Yv,X\rangle$ together with the definition
$(\curl v)^\flat=*\dd(v^\flat)$.
For \eqref{eq:curlcross}, use
$(u\times v)^\flat=*(u^\flat\wedge v^\flat)$ and the Cartan identity;
the divergence terms vanish, leaving
$\nabla_vu-\nabla_uv=-[u,v]$.
\end{proof}

\begin{lemma}[Killing fields preserve curl shells]
Let $K$ be Killing and let $v$ be divergence-free.  Then
\begin{equation}
 \curl[K,v]=[K,\curl v].
 \label{eq:curlnatural}
\end{equation}
Consequently, if $v\in E_{k,\sigma}$, then
$[K,v]\in E_{k,\sigma}$.
\end{lemma}

\begin{proof}
The flow of a Killing field consists of orientation-preserving isometries
in the identity component.  The Hodge star, exterior derivative, musical
maps, and hence curl are natural under such isometries.  Differentiating
that equivariance at the identity gives \eqref{eq:curlnatural}.  If
$\curl v=\mu_{k,\sigma}v$, then
$\curl[K,v]=\mu_{k,\sigma}[K,v]$.
\end{proof}

\begin{theorem}[Exact Killing--Beltrami polarization]
\label{thm:KBidentity}
Let
\[
 K_\tau\in\cK_\tau,\qquad
 v\in E_{k,\sigma},\qquad
 \tau,\sigma\in\{\pm1\}.
\]
Then
\begin{equation}
 \boxed{
 Q(K_\tau,v)=
 \frac{\mu_{k,\sigma}-\mu_{0,\tau}}
      {2\mu_{k,\sigma}}
 [K_\tau,v].
 }
 \label{eq:KBmu}
\end{equation}
Equivalently,
\begin{equation}
 \boxed{
 2Q(K_\tau,v)
 =
 \frac{k+2-2\tau\sigma}{k+2}[K_\tau,v].
 }
 \label{eq:KBk}
\end{equation}
In particular,
\begin{align}
 2Q(K_\sigma,v)&=\frac{k}{k+2}[K_\sigma,v],
 \label{eq:samechir}\\
 2Q(K_{-\sigma},v)&=\frac{k+4}{k+2}[K_{-\sigma},v].
 \label{eq:oppchir}
\end{align}
\end{theorem}

\begin{proof}
Set
\[
 \alpha=\mu_{0,\tau}=\frac{2\tau}{R},
 \qquad
 \beta=\mu_{k,\sigma}=\frac{\sigma(k+2)}{R}.
\]
By \eqref{eq:polarvec},
\[
 \nabla_{K_\tau}v+\nabla_vK_\tau
 =
 \gradg\langle K_\tau,v\rangle
 +(\alpha-\beta)K_\tau\times v.
\]
After Leray projection,
\begin{equation}
 Q(K_\tau,v)=\frac{\alpha-\beta}{2}
 \PP(K_\tau\times v).
 \label{eq:preKB}
\end{equation}
Since curl annihilates gradients,
\[
 \curl\PP(K_\tau\times v)
 =
 \curl(K_\tau\times v)
 =
 -[K_\tau,v].
\]
By the previous lemma,
$[K_\tau,v]$ belongs to the $\beta$-curl eigenspace.
Since $H^1(S^3)=0$, curl has no kernel on the divergence-free subspace, so
\[
 \PP(K_\tau\times v)
 =
 -\frac1\beta[K_\tau,v].
\]
Substitution into \eqref{eq:preKB} gives
\[
 Q(K_\tau,v)=\frac{\beta-\alpha}{2\beta}[K_\tau,v],
\]
which is \eqref{eq:KBmu}.  Inserting the explicit eigenvalues gives
\eqref{eq:KBk}--\eqref{eq:oppchir}.
\end{proof}

\begin{remark}
The two chiralities are dynamically inequivalent once $k\ge1$: the Killing
factor with the same curl sign acts with coefficient $k/(k+2)$, whereas the
opposite factor acts with coefficient $(k+4)/(k+2)$.  This asymmetry is what
later allows the full Killing algebra to separate the two signed shells even
though the Stokes operator assigns them the same eigenvalue.
\end{remark}



\section{Exact modal coefficients with one Killing input}

The interaction formula determines not only the target shell but also the
Hilbert--Schmidt coefficient used in Ref.~\cite{PreeceBatenin2026}.

Let
\[
 V_R:=\Vol(S^3_R)=2\pi^2R^3.
\]
Choose a pointwise orthonormal Killing frame
$X_1,X_2,X_3$ in one simple factor $\cK_\tau$; then
$\|X_a\|_{L^2}^2=V_R$ and, up to orientation,
\[
 [X_a,X_b]=\frac{2}{R}\varepsilon_{abc}X_c.
\]
Thus $K_a:=V_R^{-1/2}X_a$ is $L^2$-orthonormal in $\cK_\tau$.

For the action of $\cK_\tau$ on $E_{k,\sigma}$ define
\begin{equation}
 j_\tau(k,\sigma)
 :=
 \frac{k+1+\tau\sigma}{2}.
 \label{eq:jtausigma}
\end{equation}
This is exactly the appropriate $SU(2)_\tau$ spin from \eqref{eq:irreps}.

\begin{theorem}[Closed formula for the Killing--shell structure coefficient]
\label{thm:modalC}
Let
\[
 C_{(0,\tau),(k,\sigma)}^{(m,\rho)}
 =
 \left\|
 \Pi_{m,\rho}Q
 \big|_{\cK_\tau\times E_{k,\sigma}}
 \right\|_{\mathrm{HS}}^2.
\]
Then
\begin{equation}
 C_{(0,\tau),(k,\sigma)}^{(m,\rho)}=0
 \qquad
 \text{unless }(m,\rho)=(k,\sigma),
 \label{eq:Czero}
\end{equation}
and
\begin{equation}
\boxed{
\begin{aligned}
 C_{(0,\tau),(k,\sigma)}^{(k,\sigma)}
 &=
 \frac{d_k}{R^2V_R}
 \frac{(k+2-2\tau\sigma)^2}{(k+2)^2}\\
 &\quad\times
 j_\tau(k,\sigma)\bigl(j_\tau(k,\sigma)+1\bigr).
\end{aligned}}
\label{eq:Cexact}
\end{equation}
In particular this coefficient vanishes for all $k=0$, consistently with
$Q(\cK,\cK)=0$.
\end{theorem}

\begin{proof}
The target statement \eqref{eq:Czero} follows from
\cref{thm:KBidentity}.  Put
\[
 c_{k,\sigma,\tau}
 =
 \frac{k+2-2\tau\sigma}{2(k+2)}.
\]
Then $Q(K_a,v)=c_{k,\sigma,\tau}[K_a,v]$.
On the spin-$j=j_\tau(k,\sigma)$ representation, the normalized Casimir
identity gives
\[
 -\sum_{a=1}^3 [X_a,[X_a,\cdot]]
 =
 \frac{4}{R^2}j(j+1)\Id.
\]
The Lie-derivative action of each $X_a$ is skew-adjoint in $L^2$ because
$X_a$ is Killing.  If $\{v_i\}_{i=1}^{d_k}$ is an $L^2$-orthonormal basis
of $E_{k,\sigma}$, then
\[
 \sum_{a,i}\|[X_a,v_i]\|_2^2
 =
 \frac{4}{R^2}j(j+1)d_k.
\]
Replacing $X_a$ by $K_a=X_a/\sqrt{V_R}$ and multiplying by
$c_{k,\sigma,\tau}^2$ yields \eqref{eq:Cexact}.
\end{proof}

\begin{corollary}[Two explicit chiral values]
For $k\ge1$,
\begin{align}
 C_{(0,\sigma),(k,\sigma)}^{(k,\sigma)}
 &=
 \frac{d_k}{R^2V_R}
 \frac{k^2(k+4)}{4(k+2)},
 \label{eq:Csame}\\
 C_{(0,-\sigma),(k,\sigma)}^{(k,\sigma)}
 &=
 \frac{d_k}{R^2V_R}
 \frac{k(k+4)^2}{4(k+2)}.
 \label{eq:Copp}
\end{align}
\end{corollary}




\section{Killing momentum and shellwise skew dynamics}
\label{sec:zero-backreaction}

Samavaki and Tuomela proved conservation of the Killing component for
the deformation-tensor Navier--Stokes equation
\cite{SamavakiTuomela2020}. The identity below is its polarized bilinear
form; we include it because it couples the exact shell algebra to the full
phase space.

\begin{theorem}[Polarized Killing cancellation]
\label{thm:zeroback}
For every pair of smooth divergence-free vector fields $u,v$ on $\Sthree$,
\begin{equation}
 \boxed{
 \Pi_{\cK}Q(u,v)=0.
 }
 \label{eq:zeroback}
\end{equation}
Equivalently, every Killing field $K\in\cK$ satisfies
\begin{equation}
 \langle Q(u,v),K\rangle_{L^2}=0.
 \label{eq:zerobackweak}
\end{equation}
\end{theorem}

\begin{proof}
Since $u$ and $v$ are divergence-free and $\Sthree$ is closed,
integration by parts gives
\[
 \langle \nabla_u v,K\rangle_{L^2}
 =
 -\int_{\Sthree}\langle v,\nabla_uK\rangle\,\dd V,
\]
and similarly with $u,v$ interchanged.  Therefore
\[
 2\langle Q(u,v),K\rangle_{L^2}
 =
 -\int_{\Sthree}
 \Bigl(
 \langle v,\nabla_uK\rangle+
 \langle u,\nabla_vK\rangle
 \Bigr)\,\dd V.
\]
The integrand vanishes pointwise by the Killing equation
\[
 \langle\nabla_XK,Y\rangle+\langle\nabla_YK,X\rangle=0.
\]
Thus \eqref{eq:zerobackweak} holds for every $K\in\cK$, proving
\eqref{eq:zeroback}.
\end{proof}

\begin{corollary}[Exact evolution of Killing momentum]
\label{cor:Kmomentum}
Let $u$ solve the projected forced Navier--Stokes equation
\eqref{eq:NSproj}, and put
\[
 m(t)=\Pi_{\cK}u(t),\qquad
 v(t)=(I-\Pi_{\cK})u(t).
\]
Then
\begin{equation}
 \boxed{
 \dot m(t)=\Pi_{\cK}F(t).
 }
 \label{eq:mdot}
\end{equation}
In particular, for unforced flow the full six-dimensional Killing component
is conserved exactly.

The orthogonal component satisfies
\begin{equation}
 \boxed{
 \partial_t v+\nu Av+Q(v,v)+D(m(t))v=F_\perp,
 }
 \label{eq:globalhigh}
\end{equation}
where $F_\perp=(I-\Pi_{\cK})F$ and
\begin{equation}
 D(m)v:=2Q(m,v).
 \label{eq:Dglobaldef}
\end{equation}
If
\[
 v=\sum_{k\ge1,\ \sigma=\pm}v_{k,\sigma},
 \qquad
 m=m_++m_-,
\]
then the action is the shell-diagonal operator
\begin{equation}
\boxed{
\begin{aligned}
 D(m)v
 &=
 \sum_{k\ge1,\ \sigma=\pm}
 \Biggl[
 \frac{k}{k+2}[m_\sigma,v_{k,\sigma}]\\
 &\hspace{2.5em}
 +
 \frac{k+4}{k+2}[m_{-\sigma},v_{k,\sigma}]
 \Biggr].
\end{aligned}}
\label{eq:Dglobal}
\end{equation}
\end{corollary}

\begin{proof}
The Stokes operator is self-adjoint, $\ker A=\cK$, hence
$\Pi_{\cK}Au=0$.  Projecting \eqref{eq:NSproj} to $\cK$ and applying
\cref{thm:zeroback} gives \eqref{eq:mdot}.  Expanding
$Q(m+v,m+v)$, using $Q(m,m)=0$ and \cref{thm:zeroback}, gives
\eqref{eq:globalhigh}.  Formula \eqref{eq:Dglobal} follows shell by shell
from \cref{thm:KBidentity}.
\end{proof}

\begin{remark}[No single Killing drift represents the full operator]
The coefficients in \eqref{eq:Dglobal} depend on $(k,\sigma)$.
Consequently there is in general no single Killing field $Y_m\in\cK$ such
that $D(m)v=[Y_m,v]$ simultaneously on every signed shell.  The correct
global object is the shell-diagonal skew operator $D(m)$.

If $\Pi_{\cK}F\in L^1([0,T);\cK)$, then $m(t)$ remains bounded up to $T$.
Under that assumption any $L^\infty$ singularity must therefore occur in
$v(t)\in\cK^\perp$.
\end{remark}


\begin{corollary}[Affine Killing foliation]
\label{cor:affinefoliation}
Assume $\Pi_{\cK}F=0$.  For every $m_0\in\cK$, the affine leaf
\begin{equation}
 \boxed{
 \cH_{m_0}:=m_0+\cK^\perp
 }
 \label{eq:affineleaf}
\end{equation}
is invariant under the full projected Navier--Stokes flow.
More generally, for arbitrary forcing,
\begin{equation}
 m(t)=m_0+\int_0^t\Pi_{\cK}F(s)\,\dd s,
 \label{eq:forcedleaf}
\end{equation}
and the flow carries the initial leaf $\cH_{m_0}$ into the translated leaf
$\cH_{m(t)}$.
\end{corollary}

\begin{proof}
This is exactly \eqref{eq:mdot}.  If the Killing component of the force
vanishes, $m(t)=m_0$; otherwise its evolution is the explicit affine
translation \eqref{eq:forcedleaf}.
\end{proof}


\begin{remark}[Killing-momentum fibers are not Marsden--Weinstein reduced spaces]
\label{rem:notMW}
Define the linear Killing-momentum map
\begin{equation}
 \mathcal J:L^2_\sigma(T\Sthree)\longrightarrow\cK^*,
 \qquad
 \langle\mathcal J(u),K\rangle
 :=
 \langle u,K\rangle_{L^2}.
 \label{eq:linearJ}
\end{equation}
After the $L^2$ identification $\cK^*\simeq\cK$,
$\mathcal J(u)=\Pi_{\cK}u$, and therefore
\begin{equation}
 \mathcal J^{-1}(m)=m+\cK^\perp.
 \label{eq:Jfiber}
\end{equation}
Thus the affine leaves are exactly level sets of six conserved linear
functionals when $\Pi_{\cK}F=0$.

This statement should not be confused with Marsden--Weinstein symplectic
reduction \cite{MarsdenWeinstein1974}.  Classical reduction starts from a
symplectic Hamiltonian phase space and produces the quotient
$J^{-1}(\mu)/H_\mu$ by the stabilizer of the momentum value.  The viscous
Navier--Stokes equation is dissipative, and the affine set
$\mathcal J^{-1}(m)$ is not itself such a quotient.  In the inviscid Euler
limit the conserved Killing quantities admit the usual Hamiltonian/Noether
interpretation, but no symplectic-reduction claim is needed for
\cref{cor:affinefoliation}.
\end{remark}

\begin{theorem}[Shellwise unitary Killing precession]
\label{thm:shellwiseprecession}
For every $m\in\cK$, the operator $D(m)=2Q(m,\cdot)$ has the following
properties on the smooth divergence-free spectral core:
\begin{align}
 D(m)E_{k,\sigma}&\subset E_{k,\sigma},
 \label{eq:Dpreserves}\\
 [D(m),A]&=0,
 \qquad
 [D(m),\curl]=0,
 \label{eq:Dcommutes}\\
 D(m)^*&=-D(m).
 \label{eq:Dskew}
\end{align}
Consequently the nonautonomous propagator $U_m(t,s)$ solving
\[
 \partial_t U_m(t,s)+D(m(t))U_m(t,s)=0,
 \qquad U_m(s,s)=\Id,
\]
is unitary on $L^2$, preserves every signed curl shell, and preserves every
spectral Sobolev norm defined by a function of $A$.

If $v$ solves \eqref{eq:globalhigh}, then each signed-shell energy satisfies
\begin{equation}
\boxed{
\begin{aligned}
 \frac12\frac{\dd}{\dd t}\|v_{k,\sigma}\|_2^2
 &+\nu\kappa_k\|v_{k,\sigma}\|_2^2\\
 &+\left\langle
   \Pi_{k,\sigma}Q(v,v),v_{k,\sigma}
  \right\rangle\\
 &=
 \left\langle F_{k,\sigma},v_{k,\sigma}\right\rangle.
\end{aligned}}
\label{eq:shellenergy}
\end{equation}
In particular, the Killing component produces no transfer of energy between
signed curl shells; all inter-shell nonlinear energy transfer is carried by
$Q(v,v)$.
\end{theorem}

\begin{proof}
The shell invariance \eqref{eq:Dpreserves} is formula \eqref{eq:Dglobal}.
Since $A$ and $\curl$ are scalar on each signed shell, \eqref{eq:Dcommutes}
follows.  On each shell, \eqref{eq:Dglobal} is a real linear combination of
Lie derivatives $[K,\cdot]$ generated by Killing fields.  Such Lie
derivatives are skew-adjoint in $L^2$, proving \eqref{eq:Dskew}.  Standard
finite-shell approximation and closure then give a unitary propagator
preserving each spectral subspace and every $A$-Sobolev norm.

Project \eqref{eq:globalhigh} onto $E_{k,\sigma}$ and take the $L^2$ inner
product with $v_{k,\sigma}$.  The Killing-drift contribution vanishes by
skew-adjointness, while $Av_{k,\sigma}=\kappa_kv_{k,\sigma}$, yielding
\eqref{eq:shellenergy}.
\end{proof}

\begin{corollary}[High-mode energy balance independent of Killing precession]
\label{cor:highenergy}
For sufficiently regular solutions,
\begin{equation}
 \boxed{
 \frac12\frac{\dd}{\dd t}\|v\|_2^2
 +\nu\langle Av,v\rangle
 =
 \langle F_\perp,v\rangle.
 }
 \label{eq:highenergy}
\end{equation}
The conserved or forced Killing momentum can rotate high modes inside their
shells but contributes neither positive nor negative work to the high-mode
energy.
\end{corollary}

\begin{proof}
Sum \eqref{eq:shellenergy} over all signed shells and use the usual
Euler cancellation $\langle Q(v,v),v\rangle=0$.
\end{proof}




\section{Nonabelian shell connection and curvature}
\label{sec:holonomy}

Fix a signed shell $E_{k,\sigma}$ with $k\ge1$ and write
\[
 a_k=\frac{k}{k+2},
 \qquad
 b_k=\frac{k+4}{k+2}.
\]
For $m=m_\sigma+m_{-\sigma}\in\cK$,
\[
 D_{k,\sigma}(m)
 =
 a_k\,\ad_{m_\sigma}
 +
 b_k\,\ad_{m_{-\sigma}}
 \quad\text{on }E_{k,\sigma}.
\]

\begin{theorem}[Exact Lie-algebra representation defect]
\label{thm:representationdefect}
For $m,n\in\cK$,
\begin{align}
 [D_{k,\sigma}(m),D_{k,\sigma}(n)]
 &=
 a_k^2\ad_{[m_\sigma,n_\sigma]}
 \notag\\
 &\quad+
 b_k^2\ad_{[m_{-\sigma},n_{-\sigma}]},
 \label{eq:DcommD}\\
 \Omega_{k,\sigma}(m,n)
 &:=
 [D_{k,\sigma}(m),D_{k,\sigma}(n)]
 -
 D_{k,\sigma}([m,n])
 \notag\\
 &=
 -\frac{2k}{(k+2)^2}\ad_{[m_\sigma,n_\sigma]}
 \notag\\
 &\quad+
 \frac{2(k+4)}{(k+2)^2}
 \ad_{[m_{-\sigma},n_{-\sigma}]}.
 \label{eq:Omega}
\end{align}
Thus $m\mapsto D_{k,\sigma}(m)$ is generally not a representation of
$\mathfrak{so}(4)$.
\end{theorem}

\begin{proof}
The two simple Killing factors commute:
$[\cK_+,\cK_-]=0$, and their induced Lie-derivative actions commute on
$E_{k,\sigma}$.  Therefore
\[
\begin{aligned}
 &[a_k\ad_{m_\sigma}+b_k\ad_{m_{-\sigma}},
   a_k\ad_{n_\sigma}+b_k\ad_{n_{-\sigma}}]\\
 &\qquad=
 a_k^2\ad_{[m_\sigma,n_\sigma]}
 +
 b_k^2\ad_{[m_{-\sigma},n_{-\sigma}]}.
\end{aligned}
\]
Subtract
\[
 D_{k,\sigma}([m,n])
 =
 a_k\ad_{[m_\sigma,n_\sigma]}
 +
 b_k\ad_{[m_{-\sigma},n_{-\sigma}]},
\]
and use
$a_k(a_k-1)=-2k/(k+2)^2$ and
$b_k(b_k-1)=2(k+4)/(k+2)^2$.
\end{proof}

The defect has a precise geometric realization.  Let
$G=SO(4)$ and consider the trivial Hermitian bundle
\[
 \mathscr E_{k,\sigma}=G\times E_{k,\sigma}\longrightarrow G.
\]
For a left-invariant vector field $X^L$ generated by
$X\in\mathfrak{so}(4)\simeq\cK$, define
\begin{equation}
 \nabla^{k,\sigma}_{X^L}
 :=
 X^L + D_{k,\sigma}(X).
 \label{eq:connection}
\end{equation}

\begin{theorem}[Curvature of the shell-precession connection]
\label{thm:connectioncurvature}
The connection \eqref{eq:connection} is metric and its curvature on
left-invariant directions is exactly
\begin{equation}
 \boxed{
 R^{k,\sigma}(X^L,Y^L)
 =
 \Omega_{k,\sigma}(X,Y).
 }
 \label{eq:curvatureOmega}
\end{equation}
If $g(t)\in SO(4)$ has left-trivialized velocity
\[
 (L_{g(t)^{-1}})_*\dot g(t)=m(t),
\]
then parallel transport in $\mathscr E_{k,\sigma}$ is exactly the
Killing-precession equation
\begin{equation}
 \dot U(t)+D_{k,\sigma}(m(t))U(t)=0.
 \label{eq:paralleltransport}
\end{equation}
Consequently, if $g(t)$ is a closed loop, the shell evolution is the
nonabelian holonomy of this connection.
\end{theorem}

\begin{proof}
Metricity follows from the skew-adjointness of every
$D_{k,\sigma}(X)$.  Since the coefficients are left-invariant,
\[
\begin{aligned}
 R(X^L,Y^L)
 &=
 [\nabla_{X^L},\nabla_{Y^L}]
 -
 \nabla_{[X^L,Y^L]}\\
 &=
 [D(X),D(Y)]-D([X,Y]).
\end{aligned}
\]
which is \eqref{eq:curvatureOmega}.  The parallel-transport equation for a
connection written in left trivialization is precisely
\eqref{eq:paralleltransport}.
\end{proof}

\begin{remark}[What holonomy does and does not say]
For a sufficiently small group loop spanning directions $X,Y$, the leading
nontrivial holonomy term is proportional to the curvature
$\Omega_{k,\sigma}(X,Y)$, with the sign fixed by loop orientation and the
parallel-transport convention.  For a finite loop, the holonomy is a
path-ordered exponential; it is not in general
$\exp(\int\Omega)$.

Likewise, periodicity of the Lie-algebra-valued function $m(t)$ alone does
not make the reconstructed $SO(4)$ path $g(t)$ closed.  Thus the curvature
theorem gives an exact geometric interpretation of shell precession but not,
without an additional averaging argument, a theorem that rapid Killing
forcing suppresses nonlinear triadic resonance.
\end{remark}

\begin{theorem}[Coefficient flattening without operator-norm flattening]
\label{thm:noUVflattening}
Algebraically,
\begin{equation}
 D_{k,\sigma}(m)
 =
 \ad_m+
 \frac{2}{k+2}
 \bigl(\ad_{m_{-\sigma}}-\ad_{m_\sigma}\bigr).
 \label{eq:Dminusad}
\end{equation}
Thus the scalar coefficients approach those of the ordinary adjoint action.
However, the natural operator norms on the changing representation spaces
do not flatten.

Let $X_1^\tau,X_2^\tau,X_3^\tau$ be a pointwise-orthonormal Killing frame
in $\cK_\tau$, normalized by
\[
 [X_a^\tau,X_b^\tau]=\frac{2}{R}\varepsilon_{abc}X_c^\tau .
\]
On $E_{k,\sigma}$ the $SU(2)_\tau$ spin is
$j_\tau=(k+1+\tau\sigma)/2$, so
\begin{equation}
 \|\ad_{X_3^\tau}\|_{\mathrm{op},E_{k,\sigma}}
 =
 \frac{2j_\tau}{R}.
 \label{eq:generatoropnorm}
\end{equation}
Consequently, for a unit generator in the same chiral factor,
\begin{equation}
 \boxed{
 \|D_{k,\sigma}(X^\sigma)-\ad_{X^\sigma}\|_{\mathrm{op}}
 =
 \frac{2}{R},
 }
 \label{eq:noDflatten}
\end{equation}
for every $k\ge1$.  For a unit generator in the opposite factor,
\begin{equation}
 \|D_{k,\sigma}(X^{-\sigma})-\ad_{X^{-\sigma}}\|_{\mathrm{op}}
 =
 \frac{2k}{(k+2)R}
 \longrightarrow \frac{2}{R}.
 \label{eq:noDflattenopp}
\end{equation}

Likewise, for pointwise-orthonormal
$X_1^\tau,X_2^\tau$ in one simple factor, the curvature satisfies
\begin{align}
 \|\Omega_{k,\sigma}(X_1^\sigma,X_2^\sigma)\|_{\mathrm{op}}
 &=
 \frac{4k}{(k+2)R^2}
 \longrightarrow \frac{4}{R^2},
 \label{eq:curvnormsame}\\
 \|\Omega_{k,\sigma}(X_1^{-\sigma},X_2^{-\sigma})\|_{\mathrm{op}}
 &=
 \frac{4k(k+4)}{(k+2)^2R^2}
 \longrightarrow \frac{4}{R^2}.
 \label{eq:curvnormopp}
\end{align}
Thus the $O(k^{-1})$ curvature coefficients are exactly compensated by the
$O(k)$ growth of the representation generators.  The shell connection is
not asymptotically flat in this natural operator norm.
\end{theorem}

\begin{proof}
In the spin-$j$ representation of the Lie algebra normalized by
$[X_a,X_b]=(2/R)\varepsilon_{abc}X_c$, a Cartan generator has eigenvalues
$2im/R$, $m=-j,\ldots,j$.  Hence
\eqref{eq:generatoropnorm}.  Substitution into \eqref{eq:Dminusad} gives
\eqref{eq:noDflatten}--\eqref{eq:noDflattenopp}.

For curvature, take $X_1^\tau,X_2^\tau$ from one simple factor.  Then
$[X_1^\tau,X_2^\tau]=(2/R)X_3^\tau$.  Combining
\eqref{eq:Omega} with \eqref{eq:generatoropnorm}, using
$j_\sigma=(k+2)/2$ in the same chiral factor and
$j_{-\sigma}=k/2$ in the opposite factor, gives
\eqref{eq:curvnormsame}--\eqref{eq:curvnormopp}.
\end{proof}

\begin{remark}
A coefficient-level or suitably rescaled semiclassical curvature does decay
like $k^{-1}$.  This is distinct from operator-norm flatness on
$E_{k,\sigma}$ and should not be interpreted as rigid-body convergence of the
unscaled shell dynamics.
\end{remark}



\section{Exact Killing--Beltrami reductions}

\begin{theorem}[Killing plus one signed shell]
\label{thm:KplusE}
For every $k\ge0$ and $\sigma,\tau\in\{\pm1\}$,
\begin{equation}
 \cK_\tau\oplus E_{k,\sigma}
 \label{eq:halfKclosure}
\end{equation}
is an exact Navier--Stokes reducing space.  Likewise,
\begin{equation}
 \boxed{
 \cK\oplus E_{k,\sigma}
 =
 E_{0,+}\oplus E_{0,-}\oplus E_{k,\sigma}
 }
 \label{eq:fullKclosure}
\end{equation}
is exact reducing.
For $k\ge1$ these are genuinely multi-shell closures.
\end{theorem}

\begin{proof}
The Stokes operator preserves each summand.  By
\eqref{eq:diagonalzero},
\[
 Q(\cK_\tau,\cK_\tau)=0,
 \qquad
 Q(\cK,\cK)=0,
 \qquad
 Q(E_{k,\sigma},E_{k,\sigma})=0.
\]
By \cref{thm:KBidentity},
\[
 Q(\cK_\tau,E_{k,\sigma})\subset E_{k,\sigma},
 \qquad
 Q(\cK,E_{k,\sigma})\subset E_{k,\sigma}.
\]
Therefore the polarized quadratic term has no leakage from either space.
\end{proof}

\begin{corollary}[New zeros of the complete-shell leakage scalar]
Let
\[
\begin{aligned}
 S_{\tau;k,\sigma}
 &:=
 \{(0,\tau),(k,\sigma)\},\\
 S_{K;k,\sigma}
 &:=
 \{(0,+),(0,-),(k,\sigma)\}.
\end{aligned}
\]
Then
\[
 \mathfrak R(S_{\tau;k,\sigma})=0,
 \qquad
 \boxed{\mathfrak R(S_{K;k,\sigma})=0.}
\]
\end{corollary}



\section{Exact viscous nonstationary solutions}

Fix $k\ge1$ and $\sigma\in\{\pm1\}$.  Decompose a Killing field as
\[
 K=K_\sigma+K_{-\sigma},
 \qquad
 K_{\pm\sigma}\in\cK_{\pm\sigma}.
\]
Define
\begin{equation}
 D_{k,\sigma}(K)v
 :=
 \frac{k}{k+2}[K_\sigma,v]
 +
 \frac{k+4}{k+2}[K_{-\sigma},v].
 \label{eq:Ddef}
\end{equation}
By \cref{thm:KBidentity},
\begin{equation}
 D_{k,\sigma}(K)v=2Q(K,v).
 \label{eq:DQ}
\end{equation}
Notice that
\[
 D_{k,\sigma}(K)=
 \ad_{\widetilde K},
 \qquad
 \widetilde K=
 \frac{k}{k+2}K_\sigma+
 \frac{k+4}{k+2}K_{-\sigma},
\]
so $D_{k,\sigma}(K)$ is skew-adjoint on $L^2$.

\begin{theorem}[Triangular exact dynamics]
\label{thm:triangular}
Let
\[
 u(t)=K(t)+v(t),
 \qquad K(t)\in\cK,\quad v(t)\in E_{k,\sigma},
\]
and let the projected force be
$F=F_{\cK}+F_k$ in the same space.  Then \eqref{eq:NSproj} is exactly
equivalent to
\begin{align}
 \dot K&=F_{\cK},
 \label{eq:KODE}\\
 \dot v+\nu\kappa_k v+D_{k,\sigma}(K)v&=F_k.
 \label{eq:vODE}
\end{align}
Thus the original quadratic PDE becomes a triangular finite-dimensional
system: the Killing component is independent and the shell component is
linear once $K(t)$ is known.
\end{theorem}

\begin{proof}
Expand
\[
 Q(K+v,K+v)=Q(K,K)+2Q(K,v)+Q(v,v).
\]
The first and last terms vanish by \eqref{eq:diagonalzero}; the middle term
is \eqref{eq:DQ}.  Also $AK=0$ and $Av=\kappa_kv$.
Orthogonality of the Stokes eigenspaces gives
\eqref{eq:KODE}--\eqref{eq:vODE}.
\end{proof}

\begin{corollary}[Closed unforced formula]
For unforced data $K(0)=K_0$, $v(0)=v_0$,
\begin{equation}
 \boxed{
 u(t)=
 K_0+
 e^{-\nu k(k+4)t/R^2}
 e^{-tD_{k,\sigma}(K_0)}v_0.
 }
 \label{eq:exactsolution}
\end{equation}
Moreover,
\begin{equation}
 \boxed{
 \|v(t)\|_2=
 e^{-\nu k(k+4)t/R^2}\|v_0\|_2.
 }
 \label{eq:vdecay}
\end{equation}
In the inviscid limit the shell component evolves by a one-parameter
isometry representation and has constant $L^2$ norm.
\end{corollary}

\begin{proof}
In the unforced case $K(t)=K_0$.  Since $D_{k,\sigma}(K_0)$ is
time-independent and skew-adjoint, \eqref{eq:vODE} integrates to
\eqref{eq:exactsolution}; skew-adjointness gives \eqref{eq:vdecay}.
\end{proof}

\begin{remark}[Relation to exact Euler waves]
Recent work of Heslin and Preston constructs exact nonstationary Euler
solutions on curved manifolds, including the round three-sphere, by perturbing
rigid motions with suitable eigenfields \cite{HeslinPreston2026}. Exact
Navier--Stokes solutions generated by Beltrami structures have also been
studied in Euclidean settings \cite{Baron2025}.  Formula \eqref{eq:exactsolution}
is a viscous exact closure for the Ebin--Marsden viscosity used here.  The
present derivation is internal to the Preece--Batenin spectral data and does
not require importing the Heslin--Preston construction.
\end{remark}



\section{Rigidity of Killing-saturated reductions}

The previous theorem gives examples.  We now prove a converse that changes the
status of the complete-shell criterion in an important subclass.

\begin{theorem}[Shell completeness forced by the full Killing algebra]
\label{thm:Ksaturated}
Let $W$ be a finite-dimensional exact Navier--Stokes reducing space of smooth
divergence-free vector fields on $\Sthree$, and assume
\begin{equation}
 \cK=E_{0,+}\oplus E_{0,-}\subset W.
 \label{eq:containsK}
\end{equation}
Then there exists a finite set
$S\subset\mathbb N_0\times\{\pm\}$ containing $(0,+)$ and $(0,-)$ such that
\begin{equation}
 \boxed{
 W=\bigoplus_{(k,\sigma)\in S}E_{k,\sigma}.
 }
 \label{eq:Wshellunion}
\end{equation}
That is, every nonzero higher component is forced to be a \emph{complete
signed curl shell}.
\end{theorem}

\begin{proof}
Because $W$ is reducing, $AW\subset W$.  The operator $A$ is self-adjoint and
has the distinct eigenvalues
$\kappa_k=k(k+4)/R^2$ indexed by $k$; the $\kappa_k$ eigenspace is
\[
 F_k:=E_{k,+}\oplus E_{k,-}.
\]
Since $W$ is finite-dimensional, spectral interpolation in the finite set of
eigenvalues appearing in $W$ yields
\[
 W=\bigoplus_k W_k,\qquad W_k:=W\cap F_k.
\]
By assumption $W_0=F_0=\cK$.

Fix $k\ge1$ with $W_k\ne0$.  Put
\[
 a=\frac{k}{k+2},\qquad b=\frac{k+4}{k+2};
\]
then $a,b$ are nonzero and $a\ne b$.

Take $X\in\cKp$.  Since $\cKp\subset W$ and $W$ is quadratically closed,
the operator
\[
 D_X:=2Q(X,\cdot)
\]
preserves $W_k$.  By \cref{thm:KBidentity},
\begin{equation}
 D_X
 =
 a\,\rho_+(X)\oplus b\,\rho_-(X)
 \quad\text{on }F_k=E_{k,+}\oplus E_{k,-},
 \label{eq:DXblocks}
\end{equation}
where $\rho_\pm(X)=[X,\cdot]$ on $E_{k,\pm}$.

If $X,Y\in\cKp$ and $Z=[X,Y]$, then $[D_X,D_Y]$ also preserves $W_k$ and
\[
 [D_X,D_Y]
 =
 a^2\rho_+(Z)\oplus b^2\rho_-(Z).
\]
Since $\mathfrak{su}(2)$ is perfect, every $Z\in\cKp$ is a linear
combination of commutators.  Linear combinations of $D_Z$ and the preceding
commutator isolate
\[
 \rho_+(Z)\oplus0,
 \qquad
 0\oplus\rho_-(Z),
\]
because the coefficient matrix
$\bigl(\begin{smallmatrix}a&b\\a^2&b^2\end{smallmatrix}\bigr)$ is invertible.

Hence $W_k$ is invariant under the block-supported $\cKp$ actions.
The quadratic Casimir of the first block equals a positive scalar times the
orthogonal projector onto $E_{k,+}$ and vanishes on $E_{k,-}$; similarly for
the second block.  Therefore
\[
 W_k=(W_k\cap E_{k,+})\oplus(W_k\cap E_{k,-}).
\]
Call these intersections $L_{k,+}$ and $L_{k,-}$.

Now use the $\cKm$ part of the contained Killing algebra.  By
\cref{thm:KBidentity}, its action on each $E_{k,\sigma}$ again has a nonzero
coefficient.  Therefore each $L_{k,\sigma}$ is invariant under both simple
factors $\cKp$ and $\cKm$.  Thus $L_{k,\sigma}$ is an $SO(4)$-submodule of
$E_{k,\sigma}$.  By irreducibility \eqref{eq:irreps},
\[
 L_{k,\sigma}=0
 \quad\text{or}\quad
 L_{k,\sigma}=E_{k,\sigma}.
\]
Doing this for every $k$ proves \eqref{eq:Wshellunion}.
\end{proof}

\begin{corollary}[Complete classification of Killing-saturated finite reductions]
\label{cor:Ksaturatedclassification}
A finite-dimensional smooth subspace $W$ satisfying $\cK\subset W$ is an
exact reducing space if and only if
\[
 W=W_S=
 \bigoplus_{\alpha\in S}E_\alpha
\]
for a finite shell set $S$ containing $(0,+),(0,-)$ and satisfying
\begin{equation}
 \boxed{\mathfrak R(S)=0.}
 \label{eq:KsatR}
\end{equation}
Thus, inside the Killing-saturated category, the complete-shell leakage
criterion is both a test and an exhaustive classification.
\end{corollary}

\begin{proof}
The forward implication is \cref{thm:Ksaturated} followed by the criterion of Ref.~\cite{PreeceBatenin2026} \eqref{eq:Rcriterion}.  The reverse implication is exactly
\eqref{eq:Rcriterion}.
\end{proof}

\begin{corollary}[One-shell rigidity]
Fix $k\ge1$ and $\sigma$.  If
\[
 \cK\subset W\subset\cK\oplus E_{k,\sigma}
\]
and $W$ is exact reducing, then
\[
 W=\cK
 \quad\text{or}\quad
 W=\cK\oplus E_{k,\sigma}.
\]
\end{corollary}




\section{Casimir representation of Stokes dissipation}
\label{sec:casimir}

Let $X_a^+$, $X_a^-$, $a=1,2,3$, be the standard left- and right-invariant
global Killing frames on the round three-sphere, normalized so that each
triple is pointwise orthonormal.  Set
\[
 V_R=\Vol(\Sthree),\qquad
 K_a^\pm=V_R^{-1/2}X_a^\pm.
\]
Then the six $K_a^\pm$ are $L^2$-orthonormal in $\cK$.
Define
\[
 T_a^\pm:=Q(K_a^\pm,\cdot).
\]

\begin{theorem}[Exact quadratic Casimir identity]
\label{thm:Tcasimir}
On smooth divergence-free vector fields,
\begin{equation}
 \boxed{
 A=-2V_R\sum_{\tau=\pm}\sum_{a=1}^3(T_a^\tau)^2.
 }
 \label{eq:Tcasimir}
\end{equation}
Equivalently, for every $u\in H^1_\sigma(T\Sthree)$,
\begin{equation}
 \boxed{
 \langle Au,u\rangle_{L^2}
 =
 2V_R\sum_{\tau=\pm}\sum_{a=1}^3
 \|Q(K_a^\tau,u)\|_2^2.
 }
 \label{eq:Tcasimirform}
\end{equation}
\end{theorem}

\begin{proof}
It is enough to work on a signed shell $E_{k,\sigma}$ and then sum
orthogonally.  By \cref{thm:modalC}, the Hilbert--Schmidt trace per unit
vector in the shell is
\[
 \sum_{\tau=\pm}\frac{
 C_{(0,\tau),(k,\sigma)}^{(k,\sigma)}
 }{d_k}.
\]
Using \eqref{eq:Csame}--\eqref{eq:Copp},
\[
\begin{aligned}
 \sum_{\tau=\pm}\frac{C_\tau}{d_k}
 &=
 \frac{1}{R^2V_R}
 \left[
 \frac{k^2(k+4)}{4(k+2)}
 +
 \frac{k(k+4)^2}{4(k+2)}
 \right]\\
 &=
 \frac{k(k+4)}{2R^2V_R}
 =
 \frac{\kappa_k}{2V_R}.
\end{aligned}
\]
By $SO(4)$-equivariance the positive operator
$\sum_{\tau,a}(T_a^\tau)^*T_a^\tau$ is scalar on the irreducible
$E_{k,\sigma}$, and its scalar is therefore $\kappa_k/(2V_R)$.
Hence
\[
 A=2V_R\sum_{\tau,a}(T_a^\tau)^*T_a^\tau.
\]
Each $T_a^\tau$ is skew-adjoint: on every signed shell it is a real scalar
multiple of the Lie derivative $[K_a^\tau,\cdot]$, and Killing Lie
derivatives are skew-adjoint in $L^2$.  Thus
$(T_a^\tau)^*=-T_a^\tau$, yielding \eqref{eq:Tcasimir}.
\end{proof}

There is also a simpler relation with the ordinary isometry Casimir.
Define
\begin{equation}
 \Delta_{\mathrm{iso}}
 :=
 -\sum_{\tau=\pm}\sum_{a=1}^3
 \mathcal L_{X_a^\tau}^2.
 \label{eq:Diso}
\end{equation}

\begin{theorem}[Correct Stokes--isometry-Casimir relation]
\label{thm:stokesiso}
On divergence-free vector fields on the round three-sphere,
\begin{equation}
 \boxed{
 \Delta_{\mathrm{iso}}=2\,\curl^2,
 \qquad
 A=\frac12\Delta_{\mathrm{iso}}-\frac{4}{R^2}\Id.
 }
 \label{eq:stokesiso}
\end{equation}
If instead the $L^2$-normalized fields $K_a^\tau$ are used to define
$\Delta_K=-\sum_{\tau,a}\mathcal L_{K_a^\tau}^2$, then
\begin{equation}
 \boxed{
 A=\frac{V_R}{2}\Delta_K-\frac{4}{R^2}\Id.
 }
 \label{eq:stokesisoL2}
\end{equation}
\end{theorem}

\begin{proof}
On $E_{k,\sigma}$ the two $SU(2)$ spins are
\[
 j_+=\frac{k+1+\sigma}{2},
 \qquad
 j_-=\frac{k+1-\sigma}{2}.
\]
With the normalization of the pointwise-unit frames, the two Casimirs have
eigenvalues $4j_\pm(j_\pm+1)/R^2$.  Their sum is
\[
 \frac4{R^2}
 \left(j_+(j_++1)+j_-(j_-+1)\right)
 =
 \frac{2(k+2)^2}{R^2}.
\]
Since $\curl^2$ has eigenvalue $(k+2)^2/R^2$ and
$A$ has eigenvalue $k(k+4)/R^2=((k+2)^2-4)/R^2$, the identities follow
on every signed shell and hence spectrally on their direct sum.
\end{proof}

\begin{corollary}[Exact dissipation as six convective responses]
\label{cor:dissipation}
For $u\in H^1_\sigma(T\Sthree)$,
\begin{equation}
 \boxed{
 2\nu\|\Def u\|_2^2
 =
 \nu\langle Au,u\rangle
 =
 2\nu V_R
 \sum_{\tau=\pm}\sum_{a=1}^3
 \|Q(K_a^\tau,u)\|_2^2.
 }
 \label{eq:dissipation}
\end{equation}
Consequently, if the dissipation tends to infinity along a sequence of
times, then the maximum of these six response norms tends to infinity along
a subsequence.
\end{corollary}


\begin{remark}[What the six-response identity does not imply]
Equation \eqref{eq:dissipation} is an exact global $L^2$ identity.  It does
not by itself imply microlocal anisotropy, spatial polarization of a
singular core, or concentration along a preferred geometric axis.  The
rigorous conclusion is only the finite-sum statement in the corollary:
unbounded dissipation forces at least one of the six global response norms
to become unbounded along a subsequence.
\end{remark}

\begin{proof}
Integration by parts in
$A=-\PP(2\,\operatorname{div}\Def)$ gives
$\langle Au,u\rangle=2\|\Def u\|_2^2$.
Now apply \eqref{eq:Tcasimirform}.  The final statement follows because
there are only six nonnegative summands.
\end{proof}



\section{The entire maximal-torus sector is pressure-linear}

Reference~\cite{PreeceBatenin2026} proves that every smooth
$T^2=U(1)_\theta\times U(1)_z$ invariant divergence-free field has the form
\[
 u=\alpha(s)\partial_\theta+\beta(s)\partial_z,
 \qquad s=r^2/R^2,
\]
and that its nonlinear term is absorbed exactly into pressure.

Let $\cV_{T^2}$ denote the full infinite-dimensional space of these fields.

\begin{theorem}[Infinite-dimensional pressure-linearity]
\label{thm:T2Qzero}
For all $u,v\in\cV_{T^2}$,
\begin{equation}
 \boxed{Q(u,v)=0.}
 \label{eq:T2Qzero}
\end{equation}
Hence the projected Navier--Stokes equation is exactly linear on the entire
maximal-torus invariant sector.
\end{theorem}

\begin{proof}
The torus reduction of Ref.~\cite{PreeceBatenin2026} gives $Q(w,w)=0$ for every
$w\in\cV_{T^2}$.  Since $\cV_{T^2}$ is a vector space, polarization yields
\[
 2Q(u,v)=Q(u+v,u+v)-Q(u,u)-Q(v,v)=0.
\]
\end{proof}

Reference~\cite{PreeceBatenin2026} also diagonalizes the Stokes operator with Jacobi modes
\begin{align}
 e_n^\theta&=P_n^{(1,0)}(1-2s)\partial_\theta,\\
 e_n^z&=P_n^{(0,1)}(1-2s)\partial_z,
\end{align}
and
\begin{equation}
\begin{aligned}
 H_n&:=\operatorname{span}\{e_n^\theta,e_n^z\},\\
 A|_{H_n}&=\lambda_n\Id,
 \qquad
 \lambda_n=\frac{4n(n+2)}{R^2}.
\end{aligned}
\label{eq:T2lambda}
\end{equation}
The $\lambda_n$ are strictly increasing.

\begin{theorem}[Finite toric classification]
\label{thm:T2class}
A finite-dimensional subspace $W\subset\cV_{T^2}$ is exact reducing if and
only if
\begin{equation}
 \boxed{
 W=\bigoplus_{n\in F}L_n,
 \qquad
 L_n\subset H_n,
 }
 \label{eq:T2finiteclass}
\end{equation}
for some finite set $F\subset\mathbb N_0$ and arbitrary linear subspaces
$L_n$ of the two-dimensional eigenspaces $H_n$.
\end{theorem}

\begin{proof}
By \cref{thm:T2Qzero}, quadratic closure is automatic.  Thus $W$ is reducing
if and only if $AW\subset W$.  Since $A$ is self-adjoint and the eigenvalues
\eqref{eq:T2lambda} are distinct, every finite-dimensional $A$-invariant
subspace splits into its intersections with the eigenspaces $H_n$, giving
\eqref{eq:T2finiteclass}.  Conversely any such direct sum is $A$-invariant.
\end{proof}

\begin{corollary}[Global smooth toric dynamics]
For smooth $T^2$-invariant initial data and smooth $T^2$-invariant forcing,
the smooth solution remains smooth for every finite time for which the force
is smooth.  In the unforced case,
\begin{equation}
 \|u(t)-\Pi_{H_0}u_0\|_{H^s}
 \le
 e^{-12\nu t/R^2}
 \|u_0-\Pi_{H_0}u_0\|_{H^s}
 \label{eq:T2decay}
\end{equation}
for every spectral Sobolev index $s$.  Thus finite-time singularity is
impossible in the entire infinite-dimensional maximal-torus sector.
\end{corollary}

\begin{proof}
The sector is a pair of scalar Jacobi heat equations.  The first positive
eigenvalue is $\lambda_1=12/R^2$, while $H_0$ is the stationary
two-dimensional Killing subspace.  The spectral formula gives
\eqref{eq:T2decay}, and Duhamel's formula handles smooth forcing.
\end{proof}



\section{A homogeneous addition theorem for vector projectors}

The exact projector estimate in Ref.~\cite{PreeceBatenin2026} is not special to the
canonical band.  The correct principle is homogeneity plus irreducible
isotropy.

\begin{theorem}[Homogeneous-bundle evaluation theorem]
\label{thm:homogeneousprojector}
Let $M=G/H$ be a compact homogeneous Riemannian manifold and let
$\mathcal E\to M$ be a rank-$r$ real Euclidean homogeneous vector bundle
whose isotropy representation on a fiber is irreducible.  Let
$W\subset C^\infty(\mathcal E)$ be a finite-dimensional $G$-invariant
subspace and let $\Pi_W$ be the $L^2$-orthogonal projector.  If
$N=\dim W$ and $V=\Vol(M)$, then
\begin{align}
 \boxed{
 \|\Pi_W\|_{L^2\to L^\infty}
 =
 \sqrt{\frac{N}{rV}},
 }
 \label{eq:generalL2}\\
 \boxed{
 \|\Pi_W\|_{L^1\to L^\infty}
 =
 \frac{N}{rV}.
 }
 \label{eq:generalL1}
\end{align}
Both constants are sharp.
\end{theorem}

\begin{proof}
Let $K_W(x,y):\mathcal E_y\to\mathcal E_x$ be the projector kernel and
$E_x:W\to\mathcal E_x$ the evaluation map.  Then
\[
 K_W(x,y)=E_xE_y^*,
 \qquad
 K_W(x,x)=E_xE_x^*.
\]
At a fixed point $x$, $K_W(x,x)$ commutes with the isotropy action of $H$.
By irreducibility and self-adjointness it is a scalar:
\[
 K_W(x,x)=c\,\Id_{\mathcal E_x}.
\]
Transitivity makes $c$ independent of $x$.  Taking the operator trace,
\[
 N=\operatorname{tr}\Pi_W
 =\int_M\operatorname{tr}K_W(x,x)\dd V
 =rcV,
\]
so $c=N/(rV)$.

For $L^2\to L^\infty$,
\[
 |(\Pi_Wf)(x)|\le\|E_x\|\|f\|_2=\sqrt c\,\|f\|_2.
\]
Equality is attained by
$f=E_x^*\xi/\|E_x^*\xi\|_2$ for any unit $\xi$.

For $L^1\to L^\infty$,
\[
 \|K_W(x,y)\|_{\mathrm{op}}
 \le\|E_x\|\|E_y\|=c,
\]
giving the upper bound.  The diagonal value is exactly $c$; a vector-valued
approximate identity concentrated at $y=x$ and polarized along an eigenvector
of $K_W(x,x)$ approaches this value.  Hence the bound is sharp.
\end{proof}

\begin{corollary}[Every finite $SO(4)$-invariant vector sector on $S^3$]
\label{cor:anySO4projector}
Let $W$ be any finite-dimensional $SO(4)$-invariant subspace of smooth
divergence-free vector fields on $\Sthree$.  Then, with $N=\dim W$,
\begin{align}
 \boxed{
 \|\Pi_W\|_{2\to\infty}
 =
 \sqrt{\frac{N}{3V_R}},
 }\\
 \boxed{
 \|\Pi_W\|_{1\to\infty}
 =
 \frac{N}{3V_R}.
 }
\end{align}
\end{corollary}

\begin{proof}
Apply \cref{thm:homogeneousprojector} to the tangent bundle of the round
three-sphere.  The isotropy group is $SO(3)$ acting irreducibly on the
three-dimensional tangent space.
\end{proof}



\section{Sharp spectral uncertainty and adaptive shell sparsity}

\begin{corollary}[Sharp $L^2$ concentration-volume inequality]
\label{cor:uncertainty}
If $0\ne w\in W$ and $W$ is as in
\cref{cor:anySO4projector}, then
\begin{equation}
 \boxed{
 \frac{\|w\|_2^2}{\|w\|_\infty^2}
 \ge
 \frac{3V_R}{\dim W}.
 }
 \label{eq:L2uncertainty}
\end{equation}
The constant is sharp for every such $W$.
\end{corollary}

\begin{proof}
The inequality is \cref{cor:anySO4projector} applied to $w=\Pi_Ww$.
Sharpness follows from the reproducing-kernel extremizer
$E_x^*\xi$ used in \cref{thm:homogeneousprojector}.
\end{proof}

\begin{corollary}[Sharp uncertainty inside one signed curl shell]
For $w\in E_{k,\sigma}$,
\begin{equation}
 \boxed{
 \frac{\|w\|_2^2}{\|w\|_\infty^2}
 \ge
 \frac{3V_R}{(k+1)(k+3)}.
 }
 \label{eq:shelluncertainty}
\end{equation}
Equality is attained.  Any extremizer is still a curl eigenfield and hence
generates the exact unforced Navier--Stokes trajectory
\[
 u(t)=e^{-\nu k(k+4)t/R^2}w.
\]
\end{corollary}

Now let $S$ be an arbitrary finite set of signed shell labels and put
\[
 W_S=\bigoplus_{(k,\sigma)\in S}E_{k,\sigma},
 \qquad
 N_S=\sum_{(k,\sigma)\in S}(k+1)(k+3).
\]

\begin{theorem}[Adaptive rank barrier]
\label{thm:adaptiverank}
For any integrable divergence-free field $v$,
if
\begin{equation}
 \|\Pi_{W_S}v\|_\infty\ge\eta\|v\|_\infty,
 \qquad 0<\eta\le1,
 \label{eq:capture}
\end{equation}
then
\begin{equation}
 \boxed{
 N_S\ge
 3V_R\,\eta\,\frac{\|v\|_\infty}{\|v\|_1}.
 }
 \label{eq:adaptiverank}
\end{equation}
No contiguity or low-frequency assumption on $S$ is required.
\end{theorem}

\begin{proof}
By \cref{cor:anySO4projector},
\[
 \|\Pi_{W_S}v\|_\infty
 \le
 \frac{N_S}{3V_R}\|v\|_1.
\]
Combine this with \eqref{eq:capture} and rearrange.
\end{proof}

\begin{theorem}[Derivative-budget shell-count barrier]
\label{thm:derivativeshell}
Let $S$ be a finite set of $M$ signed curl shells and assume
$\Delta_1 v\in L^1$, where $\Delta_1=\curl^2$ on coexact one-forms.
Then
\begin{equation}
 \boxed{
 \|\Pi_{W_S}v\|_\infty
 \le
 \frac{R^2}{3V_R}
 \left[
   \sum_{(k,\sigma)\in S}
   \left(1-\frac1{(k+2)^2}\right)
 \right]
 \|\Delta_1v\|_1.
 }
 \label{eq:weightedM}
\end{equation}
Hence, if
$\|\Pi_{W_S}v\|_\infty\ge\eta\|v\|_\infty$, then
\begin{equation}
 \boxed{
 M
 \ge
 \frac{3V_R\,\eta}{R^2}
 \frac{\|v\|_\infty}{\|\Delta_1v\|_1}.
 }
 \label{eq:Mderivative}
\end{equation}
This lower bound is independent of the maximum frequency of the selected shells.
\end{theorem}

\begin{proof}
On $E_{k,\sigma}$,
\[
 \Delta_1=\frac{(k+2)^2}{R^2}\Id.
\]
Therefore
\[
 \Pi_{k,\sigma}v
 =
 \frac{R^2}{(k+2)^2}\Pi_{k,\sigma}\Delta_1v.
\]
The exact $L^1\to L^\infty$ projector norm gives
\begin{align*}
 \|\Pi_{k,\sigma}v\|_\infty
 &\le
 \frac{R^2}{(k+2)^2}
 \frac{(k+1)(k+3)}{3V_R}
 \|\Delta_1v\|_1\\
 &=
 \frac{R^2}{3V_R}
 \left(1-\frac1{(k+2)^2}\right)
 \|\Delta_1v\|_1.
\end{align*}
Summing over $S$ proves \eqref{eq:weightedM}.  Since every weight in the
sum is strictly less than one, the peak-capture condition implies
\eqref{eq:Mderivative}.
\end{proof}






\begin{theorem}[Fractional $L^1$ shell barrier]
\label{thm:fractionalbarrier}
Let $s>0$ and let $\Delta_1=\curl^2$ on divergence-free one-forms.
For any finite signed-shell set $S$,
\begin{equation}
 \boxed{
 \|\Pi_{W_S}v\|_\infty
 \le
 \frac{R^s}{3V_R}
 \left[
   \sum_{(k,\sigma)\in S}
   \frac{(k+1)(k+3)}{(k+2)^s}
 \right]
 \|\Delta_1^{s/2}v\|_1.
 }
 \label{eq:fractionalbarrier}
\end{equation}
Consequently, if
$\|\Pi_{W_S}v\|_\infty\ge\eta\|v\|_\infty$, then
\begin{equation}
 \boxed{
 \sum_{(k,\sigma)\in S}
 \frac{(k+1)(k+3)}{(k+2)^s}
 \ge
 \frac{3V_R\,\eta}{R^s}
 \frac{\|v\|_\infty}{\|\Delta_1^{s/2}v\|_1}.
 }
 \label{eq:fractionalcapture}
\end{equation}
For $s=2$ this recovers the weighted estimate
\eqref{eq:weightedM}.
\end{theorem}

\begin{proof}
On $E_{k,\sigma}$,
\[
 \Delta_1^{s/2}
 =
 \left(\frac{k+2}{R}\right)^s\Id.
\]
Hence
\[
 \Pi_{k,\sigma}v
 =
 \left(\frac{R}{k+2}\right)^s
 \Pi_{k,\sigma}\Delta_1^{s/2}v.
\]
Using the exact projector norm
\[
 \|\Pi_{k,\sigma}\|_{L^1\to L^\infty}
 =
 \frac{(k+1)(k+3)}{3V_R}
\]
and summing over $S$ gives \eqref{eq:fractionalbarrier}.
The capture inequality follows by rearrangement.
\end{proof}

\begin{corollary}[Explicit global $L^1$ Bessel-potential embedding for $s>3$]
\label{cor:globalfractional}
If $s>3$ and $\Delta_1^{s/2}v\in L^1$, then
\begin{equation}
 \boxed{
 \|v\|_\infty
 \le
 \frac{2R^s}{3V_R}
 \bigl[\zeta(s-2)-\zeta(s)\bigr]
 \|\Delta_1^{s/2}v\|_1.
 }
 \label{eq:globalfractional}
\end{equation}
Thus an $L^\infty$ blow-up sequence with bounded
$\|\Delta_1^{s/2}v\|_1$ is impossible for every $s>3$.
\end{corollary}

\begin{proof}
Put $j=k+2$.  Since
\[
 \frac{(k+1)(k+3)}{(k+2)^s}
 =
 \frac{j^2-1}{j^s}
 =
 j^{2-s}-j^{-s},
\]
the sum over both signs and all $k\ge0$ is
\[
 2\sum_{j=2}^\infty(j^{2-s}-j^{-s})
 =
 2\bigl[\zeta(s-2)-\zeta(s)\bigr],
\]
which converges exactly when $s>3$.  Apply
\eqref{eq:fractionalbarrier} to an increasing sequence of finite shell sets;
the right-hand side is summable in $L^\infty$, so the shell expansion
converges absolutely there and gives \eqref{eq:globalfractional}.
\end{proof}

\begin{remark}[Interpretation of the threshold]
The condition $s>3$ is the summability threshold of this $L^1$ shell
argument in spatial dimension three.  It is an $L^1$ Bessel-potential
statement, not the usual $L^2$ Sobolev condition $H^s\hookrightarrow L^\infty$.
No claim of endpoint optimality at $s=3$ is made.
\end{remark}


\begin{proposition}[Exact pole of the supercritical embedding constant]
\label{prop:zetapole}
Write the constant in \eqref{eq:globalfractional} as
\begin{equation}
 C_s(R)
 :=
 \frac{2R^s}{3V_R}
 \bigl[\zeta(s-2)-\zeta(s)\bigr]
 =
 \frac{R^{s-3}}{3\pi^2}
 \bigl[\zeta(s-2)-\zeta(s)\bigr].
 \label{eq:Csdef}
\end{equation}
Then
\begin{equation}
 \boxed{
 \lim_{s\downarrow3}(s-3)C_s(R)
 =
 \frac{1}{3\pi^2}.
 }
 \label{eq:Csresidue}
\end{equation}
\end{proposition}

\begin{proof}
Put $s=3+\varepsilon$.  Since
\[
 \varepsilon\,\zeta(1+\varepsilon)\longrightarrow1,
 \qquad
 \varepsilon\,\zeta(3+\varepsilon)\longrightarrow0,
 \qquad
 R^\varepsilon\longrightarrow1,
\]
formula \eqref{eq:Csresidue} follows from \eqref{eq:Csdef}.
\end{proof}

\begin{remark}[The pole is not an endpoint Trudinger theorem]
The residue \eqref{eq:Csresidue} describes the divergence of the constants
in the family of estimates for $s>3$.  It does not permit one to set
$s=3$ in \eqref{eq:globalfractional}, because both the spectral series and
the norm $\|\Delta_1^{s/2}v\|_1$ vary with $s$.  In particular no
$W^{3,1}\to\mathrm{BMO}$ embedding, no Trudinger--Moser exponential
inequality, and no sharp endpoint constant follow from the zeta-pole
calculation alone.  Such endpoint statements require an independent
harmonic-analysis theorem.
\end{remark}




\section{Exact shell-helicity balances}
\label{sec:helicity}

For $v_{k,\sigma}\in E_{k,\sigma}$ define its signed shell helicity
\begin{equation}
\begin{aligned}
 \mathscr H_{k,\sigma}(v)
 &:=
 \langle v_{k,\sigma},\curl v_{k,\sigma}\rangle\\
 &=
 \mu_{k,\sigma}\|v_{k,\sigma}\|_2^2,
 \qquad
 \mu_{k,\sigma}=\sigma\frac{k+2}{R}.
\end{aligned}
\label{eq:shellhelicity}
\end{equation}

\begin{theorem}[Killing transparency and shell-helicity balance]
\label{thm:shellhelicity}
For every shell,
\begin{equation}
 \boxed{
 \langle D(m)v_{k,\sigma},
         \curl v_{k,\sigma}\rangle=0.
 }
 \label{eq:Dhelicityzero}
\end{equation}
Along a sufficiently regular solution of the high-mode equation,
\begin{equation}
\boxed{
\begin{aligned}
 \frac12\frac{\dd}{\dd t}\mathscr H_{k,\sigma}
 &+
 \nu\kappa_k\mathscr H_{k,\sigma}\\
 &+
 \mu_{k,\sigma}
 \left\langle
   \Pi_{k,\sigma}Q(v,v),v_{k,\sigma}
 \right\rangle\\
 &=
 \mu_{k,\sigma}
 \left\langle
   F_{k,\sigma},v_{k,\sigma}
 \right\rangle .
\end{aligned}}
\label{eq:shellhelicitybalance}
\end{equation}
Thus Killing precession transfers neither energy nor helicity between signed
curl shells.
\end{theorem}

\begin{proof}
Since $D(m)$ is skew-adjoint and commutes with curl,
\[
\begin{aligned}
 \langle Dv,\curl v\rangle
 &=
 -\langle v,D\curl v\rangle\\
 &=
 -\langle v,\curl Dv\rangle\\
 &=
 -\langle\curl v,Dv\rangle.
\end{aligned}
\]
so the real scalar equals its own negative and is zero.  The balance
\eqref{eq:shellhelicitybalance} is obtained by multiplying the shell-energy
identity \eqref{eq:shellenergy} by the scalar $\mu_{k,\sigma}$.
\end{proof}

\begin{theorem}[High-mode helicity law on every affine Killing leaf]
\label{thm:totalhighhelicity}
Let
\[
 \mathscr H(v)=\langle v,\curl v\rangle.
\]
Then
\begin{equation}
 \boxed{
 \frac12\frac{\dd}{\dd t}\mathscr H(v)
 +
 \nu\langle Av,\curl v\rangle
 =
 \langle F_\perp,\curl v\rangle.
 }
 \label{eq:totalhighhelicity}
\end{equation}
In particular, for inviscid unforced dynamics on a fixed affine leaf,
$\mathscr H(v)$ is conserved exactly.
\end{theorem}

\begin{proof}
Sum \eqref{eq:shellhelicitybalance} over all signed shells.  It remains to
show that the nonlinear sum vanishes:
\[
 \langle Q(v,v),\curl v\rangle=0.
\]
Because $\curl v$ is divergence-free, the Leray projection can be removed.
Using the three-dimensional identity
\[
 \nabla_vv
 =
 \gradg\!\left(\frac12|v|^2\right)
 -
 v\times\curl v,
\]
the gradient term integrates to zero against $\curl v$, while the cross
product is pointwise orthogonal to $\curl v$.  Hence the nonlinear
contribution is zero.  The Killing contribution already vanishes by
\eqref{eq:Dhelicityzero}.
\end{proof}

\begin{proposition}[Signed helicity dissipation and absolute chiral damping]
\label{prop:helicitydissipation}
The viscous term in the high-mode helicity law is
\begin{equation}
 \boxed{
 \nu\langle Av,\curl v\rangle
 =
 \frac{\nu}{R}
 \sum_{k\ge1}\kappa_k(k+2)
 \left(
   \|v_{k,+}\|_2^2-\|v_{k,-}\|_2^2
 \right).
 }
 \label{eq:signedhelicitydiss}
\end{equation}
It is therefore sign-indefinite as a functional of a mixed-chirality field.

Nevertheless viscosity is not antidissipative on the negative-helicity
sector.  With nonlinear transfer and forcing suppressed, every shell obeys
\begin{equation}
\begin{aligned}
 \left.\frac{\dd}{\dd t}\mathscr H_{k,\sigma}\right|_{\nu}
 &=
 -2\nu\kappa_k\mathscr H_{k,\sigma},\\
 \left.\frac{\dd}{\dd t}|\mathscr H_{k,\sigma}|\right|_{\nu}
 &=
 -2\nu\kappa_k|\mathscr H_{k,\sigma}|.
\end{aligned}
\label{eq:absolutechiraldecay}
\end{equation}
Thus viscosity drives both positive and negative shell helicities toward
zero.

If
\begin{equation}
 \mathscr H_{\rm abs}(v)
 :=
 \sum_{k,\sigma}
 |\mu_{k,\sigma}|\|v_{k,\sigma}\|_2^2,
 \label{eq:absolutehelicity}
\end{equation}
then its purely viscous contribution satisfies
\begin{equation}
 \boxed{
 \left.\frac{\dd}{\dd t}\mathscr H_{\rm abs}(v)\right|_{\nu}
 =
 -2\nu
 \sum_{k,\sigma}
 \kappa_k|\mu_{k,\sigma}|\|v_{k,\sigma}\|_2^2
 \le0.
 }
 \label{eq:absolutehelicitydecay}
\end{equation}
\end{proposition}

\begin{proof}
Insert
$Av_{k,\sigma}=\kappa_kv_{k,\sigma}$ and
$\curl v_{k,\sigma}=\sigma(k+2)R^{-1}v_{k,\sigma}$
into the inner product.  Equation \eqref{eq:absolutechiraldecay} follows
from the shell balance with only the viscous term retained, and summation
gives \eqref{eq:absolutehelicitydecay}.
\end{proof}

\begin{remark}
The sign-indefiniteness of \eqref{eq:signedhelicitydiss} reflects the signed
definition of total helicity, not a viscous instability of one chirality.
Likewise, a fixed sign of an azimuthal physical-space swirl does not by
itself determine the spectral imbalance between $E_{k,+}$ and $E_{k,-}$.
No chiral singularity-selection mechanism is inferred here.
\end{remark}

\begin{remark}[No chiral selection rule is inferred]
The conservation law constrains the signed sum of nonlinear shell transfers,
but it does not by itself forbid direct interactions between
$E_{k,+}$ and $E_{\ell,-}$.  A claimed ``chiral valve'' would require an
explicit Clebsch--Gordan analysis of the polarized map
$Q(E_{k,+},E_{\ell,-})$ and its target representations.  No such additional
selection rule is assumed here.
\end{remark}




\section{Sharp limitations of stronger classification and entropy claims}
\label{sec:limitations}


\begin{theorem}[Grassmannian families of exact pressure-linear reductions]
\label{thm:grassmannreductions}
Fix a signed curl shell $E_{k,\sigma}$.  Every linear subspace
$L\subset E_{k,\sigma}$ is an exact reducing space:
\begin{equation}
 AL\subset L,
 \qquad
 Q(L,L)=0.
 \label{eq:grassmannexact}
\end{equation}
Consequently, for every $1\le d\le d_k=(k+1)(k+3)$, the whole Grassmannian
\begin{equation}
 \mathrm{Gr}(d,E_{k,\sigma})
 \label{eq:grassmannfamily}
\end{equation}
consists of $d$-dimensional exact Navier--Stokes reductions.

Likewise every subspace $H\subset\cK$ is exact, because
$A|_{\cK}=0$ and $Q(\cK,\cK)=0$.
\end{theorem}

\begin{proof}
The Stokes operator is scalar on $E_{k,\sigma}$:
$A=\kappa_k\Id$.  Moreover
$Q(E_{k,\sigma},E_{k,\sigma})=0$ by the Beltrami cancellation and
polarization.  Both assertions therefore restrict to every subspace
$L\subset E_{k,\sigma}$.  The Killing statement is identical with
$\kappa_0=0$ and the pressure-linearity of the full Killing algebra.
\end{proof}

\begin{corollary}[No finite atlas of all reductions in dimension at most eight]
\label{cor:nosmallatlas}
For every $1\le d\le8$ there are continuous families of exact
$d$-dimensional reducing spaces on $\Sthree$.  Hence no finite list of
isometry types of the form ``Killing, toric, one shell, resonant sector''
can be an exhaustive classification of all exact reductions with
$\dim W\le8$ without additional hypotheses.
\end{corollary}

\begin{proof}
Choose $k$ with $d_k\ge d$ and apply
\cref{thm:grassmannreductions}.  The Grassmannian
$\mathrm{Gr}(d,E_{k,\sigma})$ has positive dimension whenever
$0<d<d_k$; for $d=8$ one may take any $k\ge2$, for which $d_k>8$.
\end{proof}

\begin{proposition}[The nonlinear-coherence ratio is not a spectral dichotomy]
\label{prop:gammalimitation}
Whenever the denominator is nonzero, define
\begin{equation}
 \Gamma(u)
 :=
 \frac{\|Q(u,u)\|_2}
 {\|u\|_2\,\|\Def u\|_2}.
 \label{eq:Gamma}
\end{equation}
Then $\Gamma(u)=0$ for every non-Killing field in any pressure-linear sector,
including every signed curl shell and the entire maximal-torus sector.
Therefore $\Gamma\to0$ does not imply concentration in one Beltrami shell.

Indeed there exist finite sums $u_M$ in the maximal-torus sector with
$\Gamma(u_M)=0$ and with nonzero energy in $M$ distinct Stokes levels.
After equalizing the level energies, the Shannon entropy of the coarsened
Stokes-level distribution is exactly $\log M$.
\end{proposition}

\begin{proof}
If $Q(u,u)=0$ and $\Def u\ne0$, then \eqref{eq:Gamma} vanishes.
For the second statement choose nonzero toric eigenmodes from $M$ distinct
Jacobi levels and rescale them to equal $L^2$ energies.  The toric sector is
pressure-linear, so the numerator remains zero.  Orthogonality of distinct
Stokes eigenspaces gives level probabilities $1/M$ and hence entropy
$\log M$.
\end{proof}

\begin{remark}
Thus neither $\Gamma(u)\to0$ nor a lower bound
$\Gamma(u)\ge\gamma_0>0$ supplies, by itself, a theorem asserting exactly two
possible forms of spectral escape.  In particular, entropy divergence still
requires a separate energy-spreading hypothesis, and shell-count divergence
requires the derivative-budget assumptions stated in the companion note.
\end{remark}

The rigidity theorem does not classify arbitrary finite-dimensional reducing
spaces outside the Killing-saturated category. The following resonance-slope
statement records one precise limitation of a narrower graph ansatz.
\label{sec:resonance-analysis}

A proposed strengthening introduces a family of three-dimensional graph
subalgebras $H_c\subset\cK$ and, for a given wave number $k$, a cancellation
condition
\[
 c=c_*(k):=\frac{k}{k+4}.
\]
Even when such a family has been constructed in a specified ansatz, the
monotonicity of $c_*(k)$ proves only the following narrow statement.

\begin{proposition}[Single-slope incompatibility across distinct levels]
\label{prop:singleslope}
Assume a fixed graph subalgebra $H_c$ can preserve a nonzero partial
subspace in level $k$ only when $c=c_*(k)=k/(k+4)$.  Then the same fixed
$H_c$ cannot satisfy that cancellation condition simultaneously for two
distinct levels $k_1\ne k_2$.
\end{proposition}

\begin{proof}
The function
\[
 c_*(k)=\frac{k}{k+4}
\]
is strictly increasing for $k\ge0$, since
\[
 c_*'(k)=\frac4{(k+4)^2}>0.
\]
Hence $c_*(k_1)=c_*(k_2)$ implies $k_1=k_2$.
\end{proof}

\begin{remark}[Why this is not a global classification theorem]
\label{rem:notglobalclass}
The preceding proposition says nothing by itself about
$Q(L_i,L_j)$ for $i\ne j$, about graph subalgebras outside the chosen
one-parameter family, about invariant spaces that are not of graph type, or
about toric reductions that mix spectral data differently.  Therefore
monotonicity of $c_*(k)$ cannot imply that \emph{all} finite-dimensional
Navier--Stokes reducing spaces are exhausted by a listed collection.
The exhaustive statement proved here is instead
\cref{cor:Ksaturatedclassification}, and it requires the full six-dimensional
Killing algebra $\cK\subset W$.
\end{remark}



Spectral escape likewise does not force Shannon entropy of the shell-energy
distribution to diverge. The following counterexample separates pointwise
concentration from energy equipartition.
\label{sec:entropy}

The projector bounds force rank or bandwidth to grow if a finite
$SO(4)$-invariant model is required to capture a fixed fraction of an
$L^\infty$ peak.  They do \emph{not} force Shannon entropy of the
$L^2$ shell-energy distribution to diverge.

For a nonzero field $u$, define
\[
 p_{k,\sigma}(u)
 =
 \frac{\|\Pi_{k,\sigma}u\|_2^2}{\|u\|_2^2},
 \qquad
 \mathcal S_{\mathrm{spec}}(u)
 =
 -\sum_{k,\sigma}p_{k,\sigma}\log p_{k,\sigma}.
\]

\begin{theorem}[Bounded energy and unbounded peak do not force entropy growth]
\label{thm:entropycounter}
There exists a sequence of smooth divergence-free fields $u_n$ on
$\Sthree$ such that
\begin{equation}
 \|u_n\|_2=1,
 \qquad
 \|u_n\|_\infty\longrightarrow\infty,
 \qquad
 \boxed{
 \mathcal S_{\mathrm{spec}}(u_n)\longrightarrow0.
 }
 \label{eq:entropycounter}
\end{equation}
Therefore no lower bound of the form
\[
 \mathcal S_{\mathrm{spec}}(u)
 \gtrsim \log\|u\|_\infty
\]
can follow from bounded $L^2$ energy and spectral projector estimates alone.
\end{theorem}

\begin{proof}
Choose an $L^2$-unit Killing field $K$ and, for each $n\ge1$, an $L^2$-unit
extremizer $w_n\in E_{n,+}$ for the sharp projector inequality.  Then
\[
 \|w_n\|_\infty
 =
 \sqrt{\frac{(n+1)(n+3)}{3V_R}}
 \asymp n.
\]
Let $\varepsilon_n=n^{-1}$ and define
\[
 u_n=\sqrt{1-\varepsilon_n}\,K+\sqrt{\varepsilon_n}\,w_n.
\]
The two terms are $L^2$-orthogonal, so $\|u_n\|_2=1$.  By the reverse
triangle inequality,
\[
 \|u_n\|_\infty
 \ge
 \sqrt{\varepsilon_n}\|w_n\|_\infty-\|K\|_\infty
 \asymp n^{1/2}-O(1)\to\infty.
\]
The shell-energy distribution has only two nonzero probabilities,
$1-\varepsilon_n$ and $\varepsilon_n$.  Hence
\[
\begin{aligned}
 \mathcal S_{\mathrm{spec}}(u_n)
 &=
 -(1-\varepsilon_n)\log(1-\varepsilon_n)\\
 &\quad
 -\varepsilon_n\log\varepsilon_n\\
 &=
 O\!\left(\frac{\log n}{n}\right)\to0.
\end{aligned}
\]
\end{proof}

\begin{remark}
This example also shows why the adaptive rank barrier does not imply a
frequency-independent lower bound on the \emph{number} of active shells.
A single shell has rank $\asymp k^2$ and can move to arbitrarily high
frequency.  Thus concentration alone permits a tradeoff between shell count
and maximum frequency; packet-specific rates for that tradeoff require additional similarity
assumptions and are not part of the present article.
\end{remark}

A true entropy statement requires an additional energy-spreading hypothesis.

\begin{proposition}[Conditional entropy lower bound]
\label{prop:conditionalentropy}
Suppose a set $S$ of $M$ signed shells carries total energy probability
$q=\sum_{\alpha\in S}p_\alpha\ge\delta>0$ and
\[
 p_\alpha\le\frac{C}{M}
 \qquad(\alpha\in S).
\]
Then, whenever $M\ge C$,
\begin{equation}
 \boxed{
 \mathcal S_{\mathrm{spec}}(u)
 \ge
 \delta\log\frac{M}{C}.
 }
 \label{eq:conditionalentropy}
\end{equation}
\end{proposition}

\begin{proof}
For $\alpha\in S$,
$-\log p_\alpha\ge\log(M/C)$.  Therefore
\[
 \mathcal S_{\mathrm{spec}}(u)
 \ge
 \sum_{\alpha\in S}-p_\alpha\log p_\alpha
 \ge
 q\log(M/C)
 \ge
 \delta\log(M/C).
\]
\end{proof}



\section{Scope and independence from the singularity-transfer problem}
\label{sec:scope}

Every theorem in the present article is an intrinsic statement about
Navier--Stokes geometry, exact reduction, or spectral projection on the round
three-sphere.  No theorem here assumes the existence of a finite-time
singular similarity packet.

The motivating Zenodo preprint of Batenin and Preece
\cite{PreeceBatenin2026} additionally formulates a conditional transfer of a
forced singular branch to the sphere.  That transfer construction and its
packet-specialized analytic machinery are outside the scope of this article.
In particular, none of the Killing--Beltrami interaction,
connection-curvature, rigidity, Casimir, toric, projector,
Bessel-potential, energy, or helicity results depends on the inner-solver,
renormalized matching, or annular-cycle stability hypotheses of the transfer
construction.

\section{Conclusion}

The central structure on the round three-sphere is the exact chiral
Killing--Beltrami intertwining law. It organizes shellwise skew dynamics, a
canonical nonflat connection, exact viscous reductions, a converse rigidity
theorem for Killing-saturated spaces, and the Casimir representation of
Stokes dissipation.

This organization also separates the present contribution from established
results. Conservation of the Killing component belongs to the existing
geometric Navier--Stokes theory \cite{SamavakiTuomela2020}, and exact
nonstationary Euler solutions on the round three-sphere have independent
precedents \cite{HeslinPreston2026}. The distinct contribution here is the
full chiral Killing--Beltrami interaction law and the rigidity, connection,
and Casimir structures that follow from it.

Compact homogeneity adds the sharp projector theory and quantitative
spectral concentration bounds. These statements are intrinsic. Any
specialization to the relative singularity construction of
Ref.~\cite{PreeceBatenin2026} requires additional packet assumptions and is
not asserted here.

\begin{acknowledgments}
OpenAI ChatGPT was used substantively to test candidate
derivations, audit logical dependencies, organize literature, and assist with
\LaTeX{} preparation. The tool was directed through explicit mathematical
prompts and instructed to distinguish proved statements from conditional or
conjectural extensions. Retained outputs were checked against explicit
derivations, the Zenodo source preprint, and the cited literature. The human
authors remain responsible for final verification and for all scientific
content of the submitted manuscript.
\end{acknowledgments}

% APS strongly encourages a CRediT-style Author Contributions paragraph.
% Add it here after the confirmed author list is available.

\section*{Data Availability}
This is a purely mathematical work and no data were created or analyzed in
this study.


\begin{thebibliography}{99}

\bibitem{PreeceBatenin2026}
B.~Batenin and A.~Preece,
``Finite-dimensional spectral closure and a relative forced singular branch
for Navier--Stokes on the round three-sphere,''
Zenodo, Version 1.0 (2026), doi:10.5281/zenodo.22858542.

\bibitem{MarsdenWeinstein1974}
J.~E. Marsden and A.~Weinstein,
``Reduction of symplectic manifolds with symmetry,''
Rep. Math. Phys. \textbf{5}, 121--130 (1974),
doi:10.1016/0034-4877(74)90021-4.

\bibitem{EbinMarsden1970}
D.~G. Ebin and J.~Marsden,
``Groups of diffeomorphisms and the motion of an incompressible fluid,''
Ann. Math. \textbf{92}, 102--163 (1970),
doi:10.2307/1970699.

\bibitem{ChanCzubakDisconzi2017}
C.~H. Chan, M.~Czubak, and M.~M. Disconzi,
``The formulation of the Navier--Stokes equations on Riemannian manifolds,''
J. Geom. Phys. \textbf{121}, 335--346 (2017),
doi:10.1016/j.geomphys.2017.07.015.


\bibitem{SamavakiTuomela2020}
M.~Samavaki and J.~Tuomela,
``Navier--Stokes equations on Riemannian manifolds,''
J. Geom. Phys. \textbf{148}, 103543 (2020),
doi:10.1016/j.geomphys.2019.103543.

\bibitem{IkedaTaniguchi1978}
A.~Ikeda and Y.~Taniguchi,
``Spectra and eigenforms of the Laplacian on $S^n$ and $P^n(\mathbb C)$,''
Osaka J. Math. \textbf{15}, 515--546 (1978).

\bibitem{HeslinPreston2026}
P.~Heslin and S.~C. Preston,
``Exact non-stationary solutions of the Euler equations in two and three
dimensions,''
arXiv:2602.13929 (2026).

\bibitem{Baron2025}
A.~Baron,
``Exact solutions to the three-dimensional Navier--Stokes equations based on
the stream function approach,''
Phys. Fluids \textbf{37}, 017159 (2025),
doi:10.1063/5.0245978.

\end{thebibliography}



% FINAL AUTHOR ACTIONS BEFORE APS UPLOAD:
% - designate exactly one corresponding author in the APS submission server;
% - confirm the AI-use disclosure truthfully describes the authors' own verification;
% - confirm funding/CRediT statements if applicable;
% - confirm Physical Review submission history and any referee requests in cover letter.
\end{document}
